\documentclass[a4paper]{article}
\usepackage{amsfonts}
\usepackage{amsmath}
\usepackage{amssymb}
\usepackage{graphicx}     

\newcommand{\dint}{\displaystyle\int}
\newcommand{\llim}{\lim\limits}

\begin{document}
  
\noindent \large Auteur: Ayoub Jaaouan          \\
\noindent \large Studierichting: Informatica   \\ 
\noindent \large Oefeningenreeks 1C nr. 26.4   \\

\medskip

\normalsize

$ \dfrac{6\operatorname{cis} (80^\circ)}{2\operatorname{cis} (20^\circ)}  $\\

$ = \dfrac{6}{2}  \operatorname{cis}(80^\circ - 20^\circ) $ \\

$ = 3 \operatorname{cis}(20^\circ)$ \\

$ = 3 \left( \cos 60^\circ + i \sin 60^\circ \right) $ \\

$ = 3 \left( \dfrac{1}{2} + i \cdot \dfrac{\sqrt{3}}{2} \right) $ \\

$ = 3 \left( \dfrac{1}{2} + i \cdot \dfrac{\sqrt{3}}{2} \right) $ \\

$ = \dfrac{3}{2} + \dfrac{3\sqrt{3}}{2}i $ \\






\begin{thebibliography}{999}
%\bibitem{a} Referentie
\end{thebibliography}
\end{document}